% ============================================================
% はじめに：講義の概要と、このガイドの使い方
% ============================================================
\colorlet{phase}{mone}\colorlet{phasetint}{monetint}
\intro{はじめに：講義の概要とこのガイドの使い方}

\topic{この講義で学ぶこと}
「経営情報」は、情報をヒト・モノ・カネと並ぶ経営資源として捉え、企業活動と情報技術の関係を学ぶ科目です。
2026年度の本講義は、生成AIの活用を軸に据えています。AIが情報を集め、読み解き、提案する時代に、
経営に携わる人は何を判断し、何を担うのか。この問いを、毎回の演習で実際にAIを使い、その出力を検証しながら考えます。

本講義の到達目標は次の3つです。
\begin{itemize}
  \item \textbf{理解する}　経営情報の基本概念と、生成AIを含む情報技術のしくみを体系的に理解する。
  \item \textbf{読み解く}　AIを分析パートナーとして、データや事例から企業・組織の課題を捉える。
  \item \textbf{考える}　AIを経営に活かす方法と、その責任ある使い方を自分の言葉で提案する。
\end{itemize}

\topic{開講の形態}
\par\noindent\begin{gtab}{colspec={Q[l,wd=30mm] X[l]}}
  開講期間 & 2026年10月3日（土）〜 11月21日（土）　土曜・全8日／全15回 \\
  時間割 & 1限 9:00--10:30 ／ 2限 10:40--12:10（初回10月3日は1限のみ、第2回以降は1・2限の連続授業） \\
  実施形態 & オンライン（接続方法は講義用 Teams で案内） \\
  使用するAI & 大学が提供する Copilot（大学のアカウントでサインイン）。大学のAI利用規程が本ガイドの記述より優先する。 \\
  資料・課題 & 講義資料・課題・履修ガイドは講義用 Teams に掲載。課題の提出先は Teams の Class Notebook。 \\
\end{gtab}

\topic{本講義の AI Policy}
すべての回で生成AIを使う演習を行います。AIとの向き合い方として、次の3つを講義全体の方針とします。
\begin{point}{使う・検証する・区別する}
\begin{itemize}
  \item \textbf{使う}　毎回の演習で生成AIを実際に使い、経営情報の道具として身につける。
  \item \textbf{検証する}　出力を鵜呑みにせず、根拠と限界を確かめてから採用する。
  \item \textbf{区別する}　AIの寄与と自分の判断を分けて示し、判断の責任は自分が持つ。
\end{itemize}
\end{point}
AIの利用を隠す必要はありません。むしろ、どこでどうAIを使ったかを「AI活用記録」として提出物に添えることが求められます
（書き方は付録を参照）。AI活用の透明性は評価の対象です。

\topic{講義予定}
全15回を4つのモジュールで構成します。各回の「AI演習」は、本ガイドの演習1として収録しています。
\par\noindent\begin{gtabh}{colspec={Q[c,wd=9mm] Q[c,wd=17mm] X[l] X[l,wd=62mm]},row{2-4}={bg=monetint!50},row{5-8}={bg=mtwotint!50},row{9-12}={bg=mthreetint!50},row{13-16}={bg=mfourtint!50}}
  回 & 日程 & 題目 & AI演習 \\
  1 & 10/3 1限 & ガイダンス／AI時代の経営情報 & 「経営情報とは何か」をAIに尋ね、回答をどこまで信じてよいかを検討する \\
  2 & 10/10 1限 & 情報と意思決定 & 身近な意思決定の事例を、AIと対話しながら段階に分解する \\
  3 & 10/10 2限 & 生成AIのしくみと限界 & 同じ問いを異なる指示で投げ、出力の違いから使い方の勘所をつかむ \\
  4 & 10/17 1限 & 企業活動と情報システムの全体像 & 身近な企業の情報システム構成をAIで推定し整理する \\
  5 & 10/17 2限 & 業務プロセスとAIによる自動化 & 日常の業務を手順に分解し、AIに任せられる段階を分析する \\
  6 & 10/24 1限 & 組織・人材とAI協働 & AI導入に賛成・反対の立場をAIに演じさせ、論点を整理する \\
  7 & 10/24 2限 & ケーススタディ：企業のAI活用 & 複数の事例をAIで要約・比較し、共通する成功要因を討議する \\
  8 & 10/31 1限 & データの収集と整備 & 公開統計データをAIに読み込ませ、分析できる形に整える \\
  9 & 10/31 2限 & 可視化と探索的分析 & 同じデータから複数のグラフを作らせ、最も伝わる表現を選ぶ \\
  10 & 11/7 1限 & 予測と分析の考え方 & AIに予測分析を行わせ、前提と限界を説明させて妥当性を確かめる \\
  11 & 11/7 2限 & データに基づく意思決定 & 分析結果をもとに、AIと共同で1枚の意思決定メモを作成する \\
  12 & 11/14 1限 & デジタル戦略とAI戦略 & 選んだ企業のAI戦略を、AIを使って批判的に評価する \\
  13 & 11/14 2限 & AIのリスクとガバナンス & AIの出力に含まれる偏りや問題点を検出し、対策を考える \\
  14 & 11/21 1限 & 最終発表：AI活用提案 & AIの寄与と自分の判断を分けて発表する \\
  15 & 11/21 2限 & 総括：AI時代の経営情報 & 講義全体をAIと共に振り返り、学んだことを自分の言葉でまとめる \\
\end{gtabh}
日程・内容は変更になる場合があります。最新の情報は講義サイトと講義用 Teams で確認してください。

\topic{出席確認のしくみ}
出席確認は毎回の講義で行います。講義サイトの「出席確認」ページに各回のAIクイズ（4択5問）があり、
学籍番号と氏名を入力して5問すべてに正解すると、サーバー時刻に基づく提出時間と確認コードを記した
出席カードが表示されます。その画面のスクリーンショットを Teams の Class Notebook に提出してください。
\begin{itemize}
  \item 選択肢の表示順はページを開くたびに入れ替わります。本ガイドでは (a)〜(d) の順で掲載しています。
  \item 間違えた問題には解説が表示されます。本ガイドの「出席確認クイズの解説」には、正解の理由に加えて、
        誤りの選択肢がなぜ誤りかと、講義内容とのつながりを書いています。
  \item 確認コードは学籍番号・氏名・提出時間から計算され、教員側で再計算して照合します。
        学籍番号と氏名は正確に入力してください。
\end{itemize}

\topic{演習の取り組み方}
各回の演習は3つです。演習1はその回の「AI演習」で、講義中にも取り組みます。演習2・3は復習用で、
知識の確認と応用を扱います。模範解答は「唯一の正解」ではなく、答え方の水準を示す例です。
自分の答えと見比べて、何が足りないか、何が違うかを考えてください。
\begin{checks}
  \item 演習に取り組む前に、その回の解説とキーワード表を読む。
  \item AIを使う演習では、プロンプトを記録し、出力を検証したうえで自分の言葉で書き直す。
  \item 固有名詞・数値・文献は、必ず一次資料（公式サイト、統計、論文）で確かめる。
  \item AIを使ったら、AI活用記録（付録）を書く。使っていない場合は「使用なし」と書く。
\end{checks}
